% report.tex
\documentclass[12pt,a4paper]{article}
\usepackage[utf8]{inputenc}
\usepackage[T1]{fontenc}
\usepackage{times}
\usepackage{geometry}
\usepackage{setspace}
\usepackage{graphicx}
\usepackage{booktabs}
\usepackage{hyperref}
\usepackage{enumitem}
\usepackage[numbers]{natbib}
\geometry{margin=1in}
\onehalfspacing
\hypersetup{
  colorlinks=true,
  linkcolor=blue,
  citecolor=blue,
  urlcolor=blue
}

\title{Quantifying and Reducing Semantic Drift in Text-to-SQL Systems}
\author{Student Name \\ Department / Institution \\ Project Supervisor}
\date{May 2026}

\begin{document}

% Title page
\maketitle
\thispagestyle{empty}
\clearpage

% Abstract
\begin{abstract}
This report presents the design, implementation, and evaluation of a constraint-driven, multi-agent system for reducing semantic drift in natural language to SQL conversion. The project defines a business ontology, implements a multi-stage agent pipeline, and introduces a composite semantic drift metric to quantify alignment, constraint adherence, and result plausibility. The methodology and results demonstrate how the system converts natural language queries into executable retrieval plans while maintaining governance-aware validation and traceability.
\end{abstract}

\clearpage
\tableofcontents
\listoffigures
\listoftables
\clearpage

\section{Introduction}
This project addresses the challenge of semantic drift in text-to-SQL systems, where the generated query may diverge from the user's intended meaning. The objective is to build a system that can parse user queries, ground them in a formal business ontology, enforce domain constraints, and verify results. The report documents the structured design decisions, the implementation of the multi-agent pipeline, and the mechanisms used to measure and reduce drift.

\section{Literature Review}
Automated translation from natural language to database queries has been studied in a range of prior works, including semantic parsing, rule-based systems, and machine learning approaches. Existing research highlights the difficulty of preserving intent and business semantics when mapping ambiguous text to structured query languages. This project extends prior work by integrating business constraint validation and a drift metric to produce more reliable retrieval results.

\section{Methodology}
\subsection{System Architecture}
The system uses a modular architecture comprising a query API, a multi-agent pipeline, an ontology engine, a drift evaluator, and a database manager. A high-level workflow is shown in Figure~\ref{fig:architecture}.

\begin{figure}[h]
  \centering
  \fbox{\parbox[c][6cm][c]{0.85\linewidth}{\centering System architecture diagram placeholder}}
  \caption{High-level system architecture and component interaction.}
  \label{fig:architecture}
\end{figure}

\subsection{Agents and Data Flow}
The pipeline contains five agents:
\begin{enumerate}[leftmargin=*]
  \item IntentParserAgent: Extracts entities, metrics, and filters from natural language input.
  \item OntologyMapperAgent: Grounds extracted elements in the business ontology and resolves semantic variants.
  \item ConstraintValidatorAgent: Validates query elements against business rules and domain constraints.
  \item ExecutionPlannerAgent: Generates execution-ready SQL templates and determines join strategies.
  \item ResultVerifierAgent: Checks the plausibility and consistency of the query output.
\end{enumerate}
State is maintained in a typed object that travels through the pipeline, ensuring traceability and consistent input/output mapping.

\subsection{Semantic Drift Metric}
The semantic drift metric is a composite score combining three components:
\begin{itemize}[leftmargin=*]
  \item Intent alignment: measures how closely the generated query matches the extracted user intent.
  \item Constraint adherence: evaluates the degree to which business rules are satisfied.
  \item Result plausibility: assesses whether the query output is statistically and logically reasonable.
\end{itemize}
The composite score is used by the orchestrator to decide whether further refinement is required.

\subsection{Implementation Environment}
The system is implemented in Python 3.12 with the following key technologies:
\begin{itemize}[leftmargin=*]
  \item FastAPI for the REST API layer.
  \item NetworkX for business ontology representation.
  \item Pydantic for typed state management and validation.
  \item pytest for automated unit and integration tests.
  \item Docker for deployment and environment consistency.
\end{itemize}

\section{Results}
The implementation demonstrates the end-to-end conversion of natural language queries into executable SQL statements with measurable drift scores. Key capabilities include ontology-based grounding, rule validation, and iterative refinement. Example outcomes include successful detection of invalid metric combinations, correction of ambiguous entity mappings, and preservation of business semantics across query execution.

\section{Discussion}
The results indicate that the proposed architecture improves reliability over a direct one-shot translation approach. The use of a formal ontology reduces misalignment by constraining query generation to valid entities and metrics. The drift metric provides an objective measure for system performance and enables convergence checks during execution.

\section{Conclusion}
This report has described a modular, traceable system for reducing semantic drift in text-to-SQL applications. The design combines business ontology modeling, a multi-agent pipeline, and a composite drift metric to achieve improved semantic fidelity. Future work may include scaling the ontology to larger domains, integrating real vector embeddings for semantic mapping, and developing a user-facing dashboard for provenance inspection.

\section{Future Work}
Potential future enhancements include:
\begin{itemize}[leftmargin=*]
  \item Expanding the ontology with additional business entities, metrics, and temporal rules.
  \item Incorporating pre-trained language models or embeddings for deeper semantic matching.
  \item Adding a visualization layer to show drift evolution and query provenance.
  \item Evaluating the system on real-world datasets and production SQL databases.
\end{itemize}

\section{References}
\begin{thebibliography}{9}
\bibitem{DBLP:semparse} J. Author, "Semantic Parsing for SQL Generation," \textit{Journal of Data Systems}, 2023.
\bibitem{DBLP:drift} A. Researcher, "Measuring Semantic Drift in Natural Language Interfaces," \textit{Conference on Data Governance}, 2024.
\end{thebibliography}

\appendix
\section{Appendix A: Selected Code Snippets}
The following excerpt illustrates the orchestrator logic for the critic loop:
\begin{verbatim}
class SemanticGroundingOrchestrator:
    async def execute_query(self, user_query, context, eval_mode):
        state = QueryState(query_text=user_query)
        for iteration in range(self.max_iterations):
            state = await self.pipeline.run(state)
            drift, metrics = self.drift_calculator.calculate(state)
            if drift < self.drift_threshold:
                break
        return state
\end{verbatim}

\section{Appendix B: Database Schema and Ontology}
The system uses a synthetic dataset with entities such as Customers, Orders, Invoices, Products, and Transactions. Schema design emphasizes referential integrity, attribute constraints, and valid join relationships.

\end{document}
